\scanpage{120}
\margin{78}
and since $\unsure{\varphi_k}\ne0$ in $Q'$, $f_nQ'=Q$, then $\dfrac{\unsure{\Phi_0}}{\unsure{\varphi_k}},\ldots,\dfrac{\unsure{\Phi_T}}{\unsure{\varphi_k}}$, and with them $F_1$ and $F_2$ are regular in $Q'$; from the relations $F_2(Q')=F_2(Q)\ne0$ follows the regularity of $R$ in $Q'$.

Let us now show that, for large $n$, the mapping $f_n$ becomes one-to-one. The equation $\unsure{\varphi^{n\eta_e}}=0$ will give a cross-section of the variety $V^{(n)}$ by some hyperplane, not containing $V^{(n)}$. Evidently, this cross-section is the image of $s$-$1$-dimensional subvariety $\mathcal E_\varphi$ of the variety $V$ (the equation $\mathcal E_\varphi$ in $U^{(i)}$ - the set $\varphi^{(i)}=0$).

According to the induction hypothesis, there exists $m(\mathcal E_\varphi)$ such that the elements of the group \margin{\unsure{if}} $H^0(\mathcal E_\varphi,O_{\mathcal E_\varphi}[mE']_{\mathcal E_\varphi})$ will give a biregular imbedding of the variety $\mathcal E_\varphi$ in a projective space. From the exact sequence (19) we obtain that this imbedding of the variety $\mathcal E_\varphi$ is the restriction of the mapping $f_m:V\to V^{(m)}$ on the subvariety $\mathcal E_\varphi$. We obtain, in particular that the image $\mathcal E_\varphi$ by the mapping $f_m$ possesses a dimension $s$-$1$. Hence it follows that the cross-section of the variety $V^{(m)}$ by a hyperplane, not containing $V^{(m)}$, possesses a dimension $s$-$1$ and, consequently,
\[
\dim V^{(m)}=s.
\]
\mspage{119}
